\section*{Preface}
\addcontentsline{toc}{section}{Preface}
Record the same noise twice, and you will never get the same waveform. Yet its statistics can be strikingly reproducible: a fluctuating voltage may have a well-defined spectrum and a precise dependence on temperature. This raises some questions. How can a voltage with zero mean still deliver energy? Why do processes with identical voltage distributions respond so differently to averaging? And when can a single measured record reveal the statistics of an entire ensemble?

Correlation is the thread that ties random waveforms to their spectra. It determines which fluctuations survive averaging, and how a circuit reshapes their mean-square contribution. It also exposes a distinction at the heart of measurement: ensemble statistics describe a process, whereas statistics computed from one finite record are only estimates. Stationarity, ergodicity, and estimation error decide what a measured trace can, and cannot, tell us. They also explain why simply recording for longer does not guarantee a smoother or more accurate spectrum.

Thermal equilibrium gives these statistics a physical scale. The energy stored in a capacitor fixes the variance of its equilibrium voltage. The resistance connected to it sets the pace of relaxation and supplies the fluctuations that balance dissipation. From this picture follow $kT/C$ noise, resistor noise, and the available noise power $kTB$. Linear network analysis then carries noise from its physical sources to a measurable output, keeping track of the correlations that may reinforce or partially cancel one another.

The circuit viewpoint leads naturally to input-referred noise, noise factor, and the role of source impedance in RF design. Noise matching and power matching solve different optimization problems, and the noise of a cascade depends on where loss and gain are placed. Phase noise brings the story back to the time domain: fluctuations in a carrier's phase appear as spectral sidebands and timing jitter, and accumulated phase diffusion broadens the carrier line. Throughout, the aim stays the same: to trace a path from a physical source of randomness to the fluctuations that a measurement actually records.
